% q0045.tex — Combining Forces
\question{
  id        = {0045},
  title     = {Combining Forces},
  source    = {OBECON},
  year      = {2026},
  round     = {Seletiva IEO -- Phase C},
  language  = {English},
  topic     = {Finance},
  subtopic  = {NPV / IRR Profiles},
  type      = {objective},
  statement = {%
    NPV profiles for Projects~F and~G show the following: Project~F has a
    positive NPV at all discount rates (NPV declines as the rate rises).
    Project~G starts positive at a 0\% discount rate, crosses zero NPV at
    10\%, and becomes negative thereafter. At a 15\% discount rate,
    NPV(F) = +\$10M and NPV(G) = $-$\$10M. What can be concluded?
  },
  choices   = {%
    \item Only Project~F should be undertaken at any cost of capital below
          20\%.
    \item Both projects should be undertaken if the cost of capital is below
          15\%.
    \item If the projects were mutually exclusive, Project~G adoption would
          depend on the cost of capital.
    \item The joint realisation of both projects has an IRR of 15\%.
  },
  answer    = {D},
  solution  = {%
    At a 15\% discount rate: NPV(F) = +\$10M and NPV(G) = $-$\$10M, so
    NPV(F + G) = \$0. This means 15\% is exactly the IRR of the combined
    portfolio --- the rate at which undertaking both projects simultaneously
    yields zero net present value. (D) is correct.
  },
}
